\chapter{Conclusioni e sviluppi futuri}
\label{chap:conclusioni}

Il presente lavoro si proponeva tre obiettivi: fornire un quadro teorico sull'applicazione della XAI alla previsione del turnover, confrontare tre framework esplicativi --- SHAP, LIME e le spiegazioni controfattuali --- analizzandone fondamenti, proprietà e limiti, e infine illustrare, attraverso un caso studio, come ciascun framework risponda a domande differenti poste dal responsabile HR.

Il quadro teorico sviluppato nel Capitolo~\ref{chap:quadro-teorico} ha collocato la XAI all'intersezione tra l'evoluzione della People Analytics verso approcci predittivi e i vincoli di trasparenza imposti dal contesto decisionale HR e dal quadro normativo europeo. L'analisi dei tre framework ha mostrato che SHAP, LIME e i controfattuali non si pongono in alternativa ma operano su piani complementari: l'attribuzione additiva di SHAP garantisce coerenza con la decisione del modello e consente una diagnosi fondata; il surrogato lineare di LIME traduce la decisione in regole accessibili a un utente non tecnico; il controfattuale indica le condizioni minime sotto cui la predizione cambierebbe, abilitando il passaggio dalla diagnosi alla prescrizione.

Il caso studio condotto sul dataset IBM HR Employee Attrition (Capitolo~\ref{chap:caso-studio}) ha reso tangibile questa complementarità. L'analisi trasversale dei tre dipendenti selezionati ha mostrato che il valore operativo delle spiegazioni non risiede nella singola tecnica ma nella loro triangolazione: la convergenza dei tre metodi rafforza la fiducia nella segnalazione e nell'intervento proposto, mentre la divergenza --- come nel caso del falso positivo dominato da caratteristiche strutturali --- segnala un pattern problematico del modello e suggerisce cautela nell'azione. Il modello XAI-assistito si configura così come strumento di supporto alla decisione, non di sua sostituzione: la spiegazione restituisce al responsabile HR gli elementi per esercitare un giudizio informato, riducendo sia il rischio di sovraffidamento acritico sia quello di rifiuto sistematico del supporto algoritmico.

La conclusione principale del lavoro è dunque che la capacità predittiva di un modello di machine learning, nel contesto HR, acquista reale valore operativo solo quando è accompagnata da un ecosistema esplicativo capace di rispondere simultaneamente alle domande \emph{perché}, \emph{come} e \emph{cosa fare}. La XAI non risolve il problema della fiducia nel modello, ma lo trasforma da scelta binaria a valutazione graduata, fondata sulla qualità e sulla coerenza delle spiegazioni che accompagnano la predizione.

\section{Sviluppi futuri}
\label{sec:sviluppi_futuri}

I limiti discussi nel Capitolo~\ref{chap:discussione} suggeriscono tre direzioni di sviluppo che potrebbero estendere e consolidare i risultati del presente lavoro.

La prima direzione riguarda l'applicazione su dati aziendali reali. Il dataset IBM HR, pur essendo un benchmark consolidato, è di natura didattica e non riflette le specificità di un contesto organizzativo operativo. L'adozione dei tre framework esplicativi su dati provenienti da un'organizzazione concreta consentirebbe di verificare se la complementarità osservata nel caso studio si mantiene anche in presenza di distribuzioni più complesse, di un numero maggiore di osservazioni e di dinamiche di turnover legate a fattori contestuali non presenti nel benchmark.

La seconda direzione concerne l'integrazione di feature qualitative. Il caso del Dipendente~240, falso negativo con un profilo di permanenza coerente su tutte le variabili osservate, ha reso evidente che i dati strutturati da soli non esauriscono i determinanti del turnover. L'inclusione di variabili derivate da survey sul clima organizzativo, da analisi del sentiment delle comunicazioni interne o da indicatori relazionali potrebbe arricchire sia la capacità predittiva del modello sia la qualità delle spiegazioni locali, fornendo al responsabile HR leve di intervento oggi non visibili.

La terza direzione riguarda l'estensione ad altri modelli predittivi. Il caso studio ha impiegato un singolo algoritmo (XGBoost) e la letteratura ha mostrato che le spiegazioni prodotte da SHAP e LIME possono variare in funzione dell'architettura del modello sottostante. Un confronto sistematico delle spiegazioni generate su modelli differenti --- reti neurali, ensemble di diversa natura, modelli lineari regolarizzati --- permetterebbe di valutare la stabilità e la generalizzabilità delle conclusioni esplicative al variare della scelta modellistica, rafforzando o circoscrivendo la portata del triangolo diagnostico proposto.
